\chapter{Pràctica}
\label{cap:practica}
\begin{exercici}
Calcula el nombre real $a$ perquè el quocient
\[
\frac{a+i}{2-i}
\]
sigui (a) un nombre real, i (b) imaginari.
\end{exercici}

\begin{solucio}
Efectuant el quocient, obtenim
\begin{eqnarray*}
\frac{a+i}{2-i} &=&\frac{(a+i)(2+i)}{(2-i)(2+i)} \\
&=&\frac{2a+2i+ai+i^{2}}{4-i^{2}} \\
&=&\frac{2a-1}{5}+\frac{2+a}{5}i
\end{eqnarray*}

(a) Perquè el quocient sigui un nombre real la part imaginària ha de ser
zero:
\[
\begin{array}{lll}
\frac{2+a}{5}=0 & \Longrightarrow & a=-2
\end{array}
\]

(b) Perquè el quocient sigui un nombre imaginari la part real ha de ser
zero:
\[
\begin{array}{lll}
\frac{2a-1}{5}=0 & \Longrightarrow & a=\frac{1}{2}
\end{array}
\]
\end{solucio}

\begin{exercici}
Troba el complex $z$ del qual se sap que $z+\overline{z}=6$ i $z\cdot
\overline{z}=25$.
\end{exercici}

\begin{solucio}
Suposem que $z=a+bi$. Aleshores $\overline{z}=a-bi$ i, per tant, les
condicions de l’enunciat s’expressen mitjançant el sistema següent:
\[
\left\{
\begin{array}{l}
(a+bi)+(a-bi)=6 \\
(a+bi)\cdot (a-bi)=25
\end{array}
\right.
\]
és a dir,
\[
\left\{
\begin{array}{l}
2a=6 \\
a^{2}+b^{2}=25
\end{array}
\right.
\]
Després, $a=3$ i
\begin{eqnarray*}
b^{2} &=&25-9 \\
b^{2} &=&16
\end{eqnarray*}
i, per tant, $b=4$ o $b=-4$. Hi ha, ja que dues solucions: $z=3+4i$ o $z=3-4i$.
\end{solucio}

\begin{exercici}
Calcula l’expressió següent en $\mathbb{C}$
\[
\sqrt{-4}+\sqrt{9}-\sqrt{-16}+\sqrt{25}+2-i
\]
\end{exercici}

\begin{solucio}
Com que
\[
\sqrt{-4}=\sqrt{4}\cdot \sqrt{-1}=2i
\]
aleshores
\[
\sqrt{-4}+\sqrt{9}-\sqrt{-16}+\sqrt{25}+2-i=2i+3-4i+5+2-i=10-3i
\]
\end{solucio}

\begin{exercici}
Escriu l’equació de segon grau amb coeficients reals que té $
z=-2-i$ com una de les seves solucions. Comprova després el resultat.
\end{exercici}

\begin{solucio}
Com els coeficients de l’equació de segon grau són nombres
reals, les solucions d’aquesta equació són dos complexos conjugats.
Per tant, si $z=-2-i$ és solució, l’altra és $\overline{z}=-2+i$.
Com que la suma d’aquests complexos és
\[
z+\overline{z}=-4
\]
i el seu producte és
\[
z\cdot \overline{z}=4-i^{2}=5
\]
aleshores, l’equació és
\[
x^{2}+4x+5=0
\]

Resolent ara aquesta equació, obtenim
\[
\frac{-4\pm \sqrt{16-20}}{2}=\frac{-4\pm \sqrt{-4}}{2}=\frac{-4\pm 2i}{2}
=\left\{
\begin{array}{l}
-2+i \\
-2-i
\end{array}
\right.
\]
\end{solucio}

\begin{exercici}
Calcula els nombres reals $a$ i $b$ perquè
\[
\frac{3a-9i}{3a+bi}=5+4i
\]
\end{exercici}

\begin{solucio}
S’ha de complir que
\begin{eqnarray*}
3a-9i &=&(5+4i)\cdot (3a+bi) \\
3a-9i &=&15a+12ai+5bi+4bi^{2} \\
3a-9i &=&(15a-4b)+(12a+5b)i
\end{eqnarray*}
Per tant,
\[
\left\{
\begin{array}{l}
15a-4b=3a \\
12a+5b=-9
\end{array}
\right.
\]
Resolent el sistema, obtenim
\[
a=-\frac{1}{3}\text{ \ \ \ i \ \ \ }b=-1
\]
\end{solucio}

\begin{exercici}
Interpreta geomètricament les transformacions següents del pla
complex: (a) $z~\longmapsto ~-z$, (b) $z~\longmapsto ~\overline{z}$, (c) $
z~\longmapsto ~iz$.
\end{exercici}

\begin{solucio}
(a) La transformació $z~\longmapsto ~-z$ és una simetria respecte
de l’origen o un gir de $180
^{\circ}
$.\begin{center}
\includegraphics[width=2.6316in,height=1.9865in,keepaspectratio]{com-10.png}
\end{center}

(b) La conjugació $z~\longmapsto ~\overline{z}$ és una simetria
respecte de l’eix $X$.\begin{center}
\includegraphics[width=1.8299in,height=1.5385in,keepaspectratio]{com-11.png}
\end{center}

(c) La transformació $z~\longmapsto ~iz$ aplical afix $(x,y)$ de $z$
en el de $iz$ que és $(-y,x)$, ja que
\[
i(x+iy)=ix+i^{2}y=-y+ix
\]
Per tant, aquesta transformació és un gir de $90
^{\circ}
$.\begin{center}
\includegraphics[width=2.2987in,height=1.548in,keepaspectratio]{com-12.png}
\end{center}
\end{solucio}

\begin{exercici}
Representa en el pla els conjunts de nombres complexos que satisfan
cada una de les condicions següents: (a) $\left\vert z\right\vert =2$,
(b) $\left\vert z-4+i\right\vert =2$, (c) $\left\vert z\right\vert <2$, (d) $
\left\vert z-4+i\right\vert \geq 2$ i (e) $\arg z=45
^{\circ}
$.
\end{exercici}

\begin{solucio}
(a) Els afixos dels complexos $z$ que satisfan la condició que el seu mòdul sigui $2$ formen una circumferència de centre l’origen i radi $2$.

(b) Els afixos dels complexos $z$ que satisfan la condició
\[
\left\vert z-(4-i)\right\vert =2
\]
és una circumferència de centre $(4,-1)$ i radi $2$.

(c) Els afixos dels complexos $z$ que satisfan la condició que el seu mòdul sigui menor que $2$ formen un disc de centre l’origen i radi $2$.

(d) Els afixos dels complexos $z$ que satisfan la condició
\[
\left\vert z-(4-i)\right\vert \geq 2
\]
formen tot el pla amb un forat delimitat per la circumferència de centre $
(4,-1)$ i radi $2$.

(e) Els afixos dels complexos $z=x+iy$ que tenen argument $45
^{\circ}
$ és la semirecta
\[
\left\{
\begin{array}{l}
y=x \\
y>0
\end{array}
\right.
\]
\end{solucio}

\begin{exercici}
Determina els conjunts següents del pla complex:
\[
\begin{array}{cccccccc}
a) & z\overline{z}=4 &  & b) & \func{Re}\frac{1}{z}=1 &  & c) & 1<\left\vert
z-i\right\vert <2
\end{array}
\]
\end{exercici}

\begin{solucio}
(a) Si $z\overline{z}=4$, aleshores
\[
\left\vert z\right\vert =\sqrt{z\overline{z}}=2
\]
i, per tant, els afixos dels complexos satisfan la condició $z
\overline{z}=4$ formen una circumferència de centre l’origen i radi $2$.

(b) Si $z=x+iy\neq 0$, aleshores
\[
\func{Re}\left( \frac{1}{z}\right) =\frac{x}{x^{2}+y^{2}}
\]
Per tant, s’ha de complir que
\[
\frac{x}{x^{2}+y^{2}}=1
\]
és a dir,
\[
x^{2}+y^{2}-x=0
\]
Aquesta equació pot expressar-se de la manera següent
\[
(x-\frac{1}{2})^{2}+y^{2}=\frac{1}{4}
\]
i, com a conseqüència, els afixos dels complexos que satisfan $\func{Re}
\frac{1}{z}=1$ és una circumferència de centre $(1/2,0)$ i radi $1/2$.

(c) Es tracta d’una corona circular, centrada en el punt $(0,1)$ i les quals
circumferències interior i exterior tenen radis 1 i 2, respectivament.
\end{solucio}

\begin{exercici}
Si l’argument del complex $z=a+bi$ és $35
^{\circ}
$, troba l’argument de $-b+ai$, $a-bi$, $-a-bi$ i $-b-ai$. Expressa
cadascun d’ells en termes del complex $z$.
\end{exercici}

\begin{solucio}
Els arguments d’aquests complexos es dedueixen de manera immediata a partir de
la seva representació gràfica.\begin{center}
\includegraphics[width=2.5382in,height=2.3514in,keepaspectratio]{com-13.png}
\end{center}Así, tenim
\[
\begin{array}{ccc}
\text{Complex} & \text{Argument} &  \\
-b+ai & 125
^{\circ}
& iz \\
a-bi & -35
^{\circ}
& \overline{z} \\
-a-bi & 215
^{\circ}
& -z \\
-b-ai & 235
^{\circ}
& \overline{iz}
\end{array}
\]
\end{solucio}

\begin{exercici}
Si el mòdul del complex $z$ és $2$, troba el mòdul dels
complexos següents: (a) $\left\vert 2z\right\vert $, (b) $\left\vert -
\overline{z}\right\vert $, (c) $\left\vert 3/z\right\vert $ i (d) $z^{3}$.
\end{exercici}

\begin{solucio}
Per les propietats del mòdul d’un nombre complex, tenim:

(a)
\[
\left\vert 2z\right\vert =\left\vert 2\right\vert \cdot \left\vert
z\right\vert =2\cdot 2=4
\]

(b)
\[
\left\vert -\overline{z}\right\vert =\left\vert \overline{z}\right\vert
=\left\vert z\right\vert =2
\]

(c)
\[
\left\vert \frac{3}{z}\right\vert =\frac{\left\vert 3\right\vert }{
\left\vert z\right\vert }=\frac{3}{2}
\]

(d)
\[
\left\vert z^{3}\right\vert =\left\vert z\right\vert ^{3}=2^{3}=8
\]
\end{solucio}

\begin{exercici}
Si $z$ és un complex unitari d’argument $\alpha $, expressa en forma
polar i trigonomètrica: $z$, $-z$, $\overline{z}$ i $-\overline{z}$.
\end{exercici}

\begin{solucio}
\[
\begin{array}{lll}
\text{Complex} & \text{Polar} & \text{Trigonomètrica} \\
z & 1_{\alpha } & \cos \alpha +i\sin \alpha \\
-z & 1_{\alpha +\pi } & -\cos \alpha -i\sin \alpha \\
\overline{z} & 1_{-\alpha } & \cos \alpha -i\sin \alpha \\
-\overline{z} & 1_{-\alpha +\pi } & -\cos \alpha +i\sin \alpha
\end{array}
\]
\end{solucio}

\begin{exercici}
Efectúa les operacions següents, expressant el resultat en forma binòmica:

\begin{enumerate}
\item[a)]
\[
1_{0}+2_{\frac{\pi }{2}}+3_{\pi }+4_{\frac{3\pi }{2}}
\]

\item[b)]
\[
\frac{4_{-\frac{\pi }{6}}\cdot (\cos 28
^{\circ}
+i\sin 28
^{\circ}
)}{(\cos 31
^{\circ}
-i\sin 31
^{\circ}
)^{2}}
\]

\item[c)]
\[
\frac{5\cdot \left( \sqrt{2}\right) _{\frac{\pi }{4}}\cdot \left( \sqrt{12}
+2i\right) \cdot (-\cos 25
^{\circ}
-i\sin 25
^{\circ}
)}{2(\cos 50
^{\circ}
-i\sin 50
^{\circ}
)}
\]
\end{enumerate}
\end{exercici}

\begin{solucio}
(a) Com que els nombres complexos no se poden sumar en forma polar,
passem cadascun d’ells a forma binòmica, i resulta
\begin{eqnarray*}
1_{0}+2_{\frac{\pi }{2}}+3_{\pi }+4_{\frac{3\pi }{2}} &=&1+2i-3-4i \\
&=&-2-2i
\end{eqnarray*}

(b) Efectuarem l’operació
\[
\frac{4_{-\frac{\pi }{6}}\cdot (\cos 28
^{\circ}
+i\sin 28
^{\circ}
)}{(\cos 31
^{\circ}
-i\sin 31
^{\circ}
)^{2}}
\]
passant tots els complexos a la forma polar. D’aquesta manera, tenim
\begin{eqnarray*}
\frac{4_{-30
^{\circ}
}\cdot 1_{28
^{\circ}
}}{(1_{-31
^{\circ}
})^{2}} &=&\frac{\left( 4\cdot 1\right) _{-30
^{\circ}
+28
^{\circ}
}}{\left( 1^{2}\right) _{2\cdot (-31
^{\circ}
)}} \\
&=&\frac{4_{-2
^{\circ}
}}{1_{-62
^{\circ}
}} \\
&=&\left( \frac{4}{1}\right) _{-2
^{\circ}
-(-62
^{\circ}
)} \\
&=&4_{60
^{\circ}
}
\end{eqnarray*}
Passant aquest resultat a forma binòmica, obtenim
\[
4(\cos 60
^{\circ}
+i\sin 60
^{\circ}
)=4\left( \frac{1}{2}+\frac{\sqrt{3}}{2}i\right) =2+2\sqrt{3}i
\]

(c) Efectuarem l’operació
\[
\frac{5\cdot \left( \sqrt{2}\right) _{\frac{\pi }{4}}\cdot \left( \sqrt{12}
+2i\right) \cdot (-\cos 25
^{\circ}
-i\sin 25
^{\circ}
)}{2(\cos 50
^{\circ}
-i\sin 50
^{\circ}
)}
\]
passant tots els complexos a la forma polar. D’aquesta manera, tenim
\begin{eqnarray*}
\frac{5_{0
^{\circ}
}\cdot \left( \sqrt{2}\right) _{45
^{\circ}
}\cdot 4_{30
^{\circ}
}\cdot 1_{25
^{\circ}
+180
^{\circ}
}}{2_{-50
^{\circ}
}} &=&\frac{\left( 20\sqrt{2}\right) _{280
^{\circ}
}}{2_{-50
^{\circ}
}} \\
&=&\left( 10\sqrt{2}\right) _{330
^{\circ}
}
\end{eqnarray*}
Passant aquest resultat a forma binòmica, obtenim
\begin{eqnarray*}
10\sqrt{2}(\cos 330
^{\circ}
+i\sin 330
^{\circ}
) &=&10\sqrt{2}\left( \frac{\sqrt{3}}{2}-\frac{1}{2}i\right) \\
&=&5\sqrt{6}-5\sqrt{2}i
\end{eqnarray*}
\end{solucio}

\begin{exercici}
Efectúa les operacions següents, expressant el resultat en forma binòmica:

\begin{enumerate}
\item[a)]
\[
(i^{5}+i^{8})^{10}
\]

\item[b)]
\[
\left( \frac{1+\sqrt{3}i}{1-i}\right) ^{10}
\]

\item[c)]
\[
\frac{(-1+i)^{5}}{(-2-2\sqrt{3}i)^{6}}
\]

\item[d)]
\[
\left[ 2(\cos 40
^{\circ}
+i\sin 40
^{\circ}
)\right] ^{3}+\left[ 2(\cos 20
^{\circ}
+i\sin 20
^{\circ}
)\right] ^{3}
\]
\end{enumerate}
\end{exercici}

\begin{solucio}
(a) Com que $i^{5}=i$ i $i^{8}=1$, hem de calcular
\[
(1+i)^{10}
\]
Per això, passarem a forma polar el complex $1+i$. El seu mòdul és
\[
r=\sqrt{1^{2}+1^{2}}=\sqrt{2}
\]
i, com el seu afix és en el primer quadrant, el seu argument és
\[
\begin{array}{lll}
\tan \alpha =1 & \Longrightarrow & \alpha =45
^{\circ}
\end{array}
\]
Després,
\[
\left[ \left( \sqrt{2}\right) _{45
^{\circ}
}\right] ^{10}=\left( \sqrt{2}\right) _{10\cdot 45
^{\circ}
}^{10}=\left( 2^{5}\right) _{90
^{\circ}
+360
^{\circ}
}=32_{90
^{\circ}
}
\]
Per tant
\[
(1+i)^{10}=32i
\]

(b) Efectuarem l’operació
\[
\left( \frac{1+\sqrt{3}i}{1-i}\right) ^{10}
\]
en forma polar. El mòdul de $1+\sqrt{3}i$ és
\[
r=\sqrt{1^{2}+\left( \sqrt{3}\right) ^{2}}=2
\]
i, com el seu afix es troba en el primer quadrant, el seu argument és
\[
\begin{array}{lll}
\tan \alpha =\sqrt{3} & \Longrightarrow & \alpha =60
^{\circ}
\end{array}
\]
El mòdul de $1-i$ és
\[
r^{\prime }=\sqrt{1^{2}+(-1)^{2}}=\sqrt{2}
\]
i, com el seu afix es troba en el quart quadrant, el seu argument és
\[
\begin{array}{lll}
\tan \alpha ^{\prime }=-1 & \Longrightarrow & \alpha ^{\prime }=-45
^{\circ}
\end{array}
\]
Per tant, l’operació que volem calcular en forma polar és
\begin{eqnarray*}
\left[ \frac{2_{60
^{\circ}
}}{\left( \sqrt{2}\right) _{-45
^{\circ}
}}\right] ^{10} &=&\left[ \left( \frac{2}{\sqrt{2}}\right) _{105
^{\circ}
}\right] ^{10} \\
&=&\left[ \left( \sqrt{2}\right) ^{10}\right] _{10\cdot 105
^{\circ}
} \\
&=&\left( 2^{5}\right) _{330
^{\circ}
+2\cdot 360
^{\circ}
} \\
&=&32_{330
^{\circ}
}
\end{eqnarray*}
Per tant,
\begin{eqnarray*}
\left( \frac{1+\sqrt{3}i}{1-i}\right) ^{10} &=&32(\cos 330
^{\circ}
+i\sin 330
^{\circ}
) \\
&=&32\left( \frac{\sqrt{3}}{2}-\frac{1}{2}i\right) \\
&=&16\sqrt{3}-16i
\end{eqnarray*}

(c) Efectuarem l’operació
\[
\frac{(-1+i)^{5}}{(-2-2\sqrt{3}i)^{6}}
\]
en forma polar. El mòdul de $-1+i$ és
\[
r=\sqrt{(-1)^{2}+1^{2}}=\sqrt{2}
\]
i, com el seu afix es troba en el segon quadrant, el seu argument és
\[
\begin{array}{lll}
\tan \alpha =-1 & \Longrightarrow & \alpha =135
^{\circ}
\end{array}
\]
El mòdul de $-2-2\sqrt{3}i$ és
\[
r^{\prime }=\sqrt{(-2)^{2}+\left( -2\sqrt{3}\right) ^{2}}=4
\]
i, com el seu afix es troba en el tercer quadrant, el seu argument és
\[
\begin{array}{lll}
\tan \alpha ^{\prime }=\sqrt{3} & \Longrightarrow & \alpha ^{\prime }=240
^{\circ}
\end{array}
\]
Per tant, l’operació que volem calcular en forma polar és
\begin{eqnarray*}
\frac{\left[ \left( \sqrt{2}\right) _{135}\right] ^{5}}{\left( 4_{240
^{\circ}
}\right) ^{6}} &=&\frac{\left[ \left( \sqrt{2}\right) ^{5}\right] _{5\cdot
135
^{\circ}
}}{\left( 4^{6}\right) _{6\cdot 240
^{\circ}
}} \\
&=&\left( \frac{4\sqrt{2}}{4096}\right) _{675
^{\circ}
-1440
^{\circ}
} \\
&=&\left( \frac{\sqrt{2}}{1024}\right) _{-45
^{\circ}
-2\cdot 360
^{\circ}
} \\
&=&\left( \frac{\sqrt{2}}{1024}\right) _{-45
^{\circ}
}
\end{eqnarray*}
Per tant,
\begin{eqnarray*}
\frac{(-1+i)^{5}}{(-2-2\sqrt{3}i)^{6}} &=&\frac{\sqrt{2}}{1024}(\cos (-45
^{\circ}
)+i\sin (-45
^{\circ}
)) \\
&=&\frac{\sqrt{2}}{1024}\left( \frac{\sqrt{2}}{2}-\frac{\sqrt{2}}{2}i\right)
\\
&=&\frac{1}{1024}-\frac{1}{1024}i
\end{eqnarray*}

(d) Passant l’expressió
\[
\left[ 2(\cos 40
^{\circ}
+i\sin 40
^{\circ}
)\right] ^{3}+\left[ 2(\cos 20
^{\circ}
+i\sin 20
^{\circ}
)\right] ^{3}
\]
a forma polar, obtenim
\[
\left( 2_{40
^{\circ}
}\right) ^{3}+\left( 2_{20
^{\circ}
}\right) ^{3}=8_{120
^{\circ}
}+8_{60
^{\circ}
}
\]
Per calcular aquesta suma de nombres complexos, passem a forma binòmica cadascun dels sumands. La forma binòmica de $8_{120
^{\circ}
}$ és
\begin{eqnarray*}
8(\cos 120
^{\circ}
+i\sin 120
^{\circ}
) &=&8\left( -\frac{1}{2}+\frac{\sqrt{3}}{2}i\right) \\
&=&-4+4\sqrt{3}i
\end{eqnarray*}
i la de $8_{60
^{\circ}
}$ és
\begin{eqnarray*}
8(\cos 60
^{\circ}
+i\sin 60
^{\circ}
) &=&8\left( \frac{1}{2}+\frac{\sqrt{3}}{2}i\right) \\
&=&4+4\sqrt{3}i
\end{eqnarray*}
Per tant,
\[
\left[ 2(\cos 40
^{\circ}
+i\sin 40
^{\circ}
)\right] ^{3}+\left[ 2(\cos 20
^{\circ}
+i\sin 20
^{\circ}
)\right] ^{3}=8\sqrt{3}i
\]
\end{solucio}

\begin{exercici}
Demostra les següents fórmules d’addició d’angles
\[
\sin (\alpha +\beta )=\sin \alpha \cdot \cos \beta +\sin \beta \cdot \cos
\alpha
\]
i
\[
\cos (\alpha +\beta )=\cos \alpha \cdot \cos \beta -\sin \alpha \cdot \sin
\beta
\]
\end{exercici}

\begin{solucio}
Com que
\[
1_{\alpha }\cdot 1_{\beta }=1_{\alpha +\beta }
\]
deduïm que
\[
(\cos \alpha +i\sin \alpha )(\cos \beta +i\sin \beta )=\cos (\alpha +\beta
)+i\sin (\alpha +\beta )
\]
Ara bé, com
\begin{eqnarray*}
(\cos \alpha +i\sin \alpha )(\cos \beta +i\sin \beta ) &=&\cos \alpha \cdot
\cos \beta +i\sin \alpha \cdot \cos \beta +i\cos \alpha \cdot \sin \beta
+i^{2}\sin \alpha \cdot \sin \beta \\
&=&(\cos \alpha \cdot \cos \beta -\sin \alpha \cdot \sin \beta )+(\sin
\alpha \cdot \cos \beta +\sin \beta \cdot \cos \alpha )i
\end{eqnarray*}
aleshores
\[
(\cos \alpha \cdot \cos \beta -\sin \alpha \cdot \sin \beta )+(\sin \alpha
\cdot \cos \beta +\sin \beta \cdot \cos \alpha )i=\cos (\alpha +\beta
)+i\sin (\alpha +\beta )
\]
i, per tant,
\[
\cos (\alpha +\beta )=\cos \alpha \cdot \cos \beta -\sin \alpha \cdot \sin
\beta
\]
i
\[
\sin (\alpha +\beta )=\sin \alpha \cdot \cos \beta +\sin \beta \cdot \cos
\alpha
\]
\end{solucio}

\begin{exercici}
Calcula $\beta $ perquè
\[
\frac{10_{3\beta }}{2_{\beta }}=\frac{5}{\sqrt{2}}+\frac{5}{\sqrt{2}}i
\]
\end{exercici}

\begin{solucio}
Primer passarem a forma polar el complex
\[
\frac{5}{\sqrt{2}}+\frac{5}{\sqrt{2}}i
\]
El seu mòdul és
\[
r=\sqrt{\left( \frac{5}{\sqrt{2}}\right) ^{2}+\left( \frac{5}{\sqrt{2}}
\right) ^{2}}=5
\]
i, com el seu afix és en el primer quadrant, el seu argument és
\[
\begin{array}{lll}
\tan \alpha =1 & \Longrightarrow & \alpha =45
^{\circ}
\end{array}
\]
Aleshores podem escriure
\[
\frac{10_{3\beta }}{2_{\beta }}=5_{45
^{\circ}
}
\]
és a dir,
\[
5_{2\beta }=5_{45
^{\circ}
}
\]
Per tant,
\[
2\beta =45
^{\circ}
+k\cdot 360
^{\circ}
\]
o sigui,
\[
\beta =22.5
^{\circ}
+k\cdot 180
^{\circ}
\]
\end{solucio}

\begin{exercici}
A partir de la fórmula de Moivre expressa $\sin 3\alpha $ i $\cos 3\alpha
$ en termes de $\sin \alpha $ i $\cos \alpha $.
\end{exercici}

\begin{solucio}
Per la fórmula de Moivre, tenim
\[
\left( \cos \alpha +i\sin \alpha \right) ^{3}=\cos 3\alpha +i\sin 3\alpha
\]
Desenvolupant $\left( \cos \alpha +i\sin \alpha \right) ^{3}$ per la fórmula del binomi de Newton, obtenim
\begin{eqnarray*}
\left( \cos \alpha +i\sin \alpha \right) ^{3} &=&\binom{3}{0}\cos ^{3}\alpha
+\binom{3}{1}\cos ^{2}\alpha \cdot i\sin \alpha +\binom{3}{2}\cos \alpha
\cdot i^{2}\sin ^{2}\alpha +\binom{3}{3}i^{3}\sin ^{3}\alpha \\
&=&\binom{3}{0}\cos ^{3}\alpha +\binom{3}{1}\cos ^{2}\alpha \cdot i\sin
\alpha -\binom{3}{2}\cos \alpha \cdot \sin ^{2}\alpha -\binom{3}{3}i\sin
^{3}\alpha
\end{eqnarray*}
Mitjançant el triangle de Tartaglia calculem els nombres
combinatoris:
\[
\begin{array}{lllllll}
&  & 1 &  & 1 &  &  \\
& 1 &  & 2 &  & 1 &  \\
1 &  & 3 &  & 3 &  & 1
\end{array}
\]
Després,
\begin{eqnarray*}
\left( \cos \alpha +i\sin \alpha \right) ^{3} &=&\cos ^{3}\alpha +3\cos
^{2}\alpha \cdot i\sin \alpha -3\cos \alpha \cdot \sin ^{2}\alpha -i\sin
^{3}\alpha \\
&=&(\cos ^{3}\alpha -3\cos \alpha \cdot \sin ^{2}\alpha )+(3\cos ^{2}\alpha
\cdot \sin \alpha -\sin ^{3}\alpha )i
\end{eqnarray*}
Per tant,
\[
(\cos ^{3}\alpha -3\cos \alpha \cdot \sin ^{2}\alpha )+(3\cos ^{2}\alpha
\cdot \sin \alpha -\sin ^{3}\alpha )i=\cos 3\alpha +i\sin 3\alpha
\]
i, com a conseqüència, deduïm
\[
\cos 3\alpha =\cos ^{3}\alpha -3\cos \alpha \cdot \sin ^{2}\alpha
\]
i
\[
\sin 3\alpha =3\cos ^{2}\alpha \cdot \sin \alpha -\sin ^{3}\alpha
\]
\end{solucio}

\begin{exercici}
Calcula les arrels sisenes de $-1$ i els afixos corresponents.
Quina figura formen aquests punts? Calcula la suma de
aquestes arrels.
\end{exercici}

\begin{solucio}
El complex $-1$ en forma polar és
\[
1_{180
^{\circ}
}
\]
Per tant,
\[
\sqrt[6]{1_{180
^{\circ}
}}=1_{\frac{180
^{\circ}
}{6}+k\cdot \frac{360
^{\circ}
}{6}}=\left\{
\begin{array}{ll}
1_{30
^{\circ}
} & k=0 \\
1_{90
^{\circ}
} & k=1 \\
1_{150
^{\circ}
} & k=2 \\
1_{210
^{\circ}
} & k=3 \\
1_{270
^{\circ}
} & k=4 \\
1_{330
^{\circ}
} & k=5
\end{array}
\right.
\]
Passant aquests complexos a forma binòmica s’obtindran les
coordenades dels afixos de les arrels. Así, tenim
\[
\begin{array}{llll}
\text{Polar} & \text{Trigonomètrica} & \text{Binòmica} & \text{Afix}
\\
1_{30
^{\circ}
} & \cos 30
^{\circ}
+i\sin 30
^{\circ}
& \frac{\sqrt{3}}{2}+\frac{1}{2}i & \left( \frac{\sqrt{3}}{2},\frac{1}{2}
\right) \\
1_{90
^{\circ}
} & \cos 90
^{\circ}
+i\sin 90
^{\circ}
& i & (0,1) \\
1_{150
^{\circ}
} & \cos 150
^{\circ}
+i\sin 150
^{\circ}
& -\frac{\sqrt{3}}{2}+\frac{1}{2}i & \left( -\frac{\sqrt{3}}{2},\frac{1}{2}
\right) \\
1_{210
^{\circ}
} & \cos 210
^{\circ}
+i\sin 210
^{\circ}
& -\frac{\sqrt{3}}{2}-\frac{1}{2}i & \left( -\frac{\sqrt{3}}{2},-\frac{1}{2}
\right) \\
1_{270
^{\circ}
} & \cos 270
^{\circ}
+i\sin 270
^{\circ}
& -i & (0,-1) \\
1_{330
^{\circ}
} & \cos 330
^{\circ}
+i\sin 330
^{\circ}
& \frac{\sqrt{3}}{2}-\frac{1}{2}i & \left( \frac{\sqrt{3}}{2},-\frac{1}{2}
\right)
\end{array}
\]
Els afixos de les arrels sisenes de $-1$ formen un hexàgon regular
inscrit en una circumferència de centre l’origen i radi $1$.\begin{center}
\includegraphics[width=2.476in,height=2.476in,keepaspectratio]{com-7.png}
\end{center}És
evident que la suma de aquestes arrels és zero, ja que és zero la suma de
els seus vectors associats.
\end{solucio}

\begin{exercici}
Calcula la arrel quadrada de $3+4i$ sense utilitzar la forma polar.
\end{exercici}

\begin{solucio}
Si $a+bi$ és una arrel quadrada de $3+4i$, aleshores
\[
(a+bi)^{2}=3+4i
\]
Després,
\begin{eqnarray*}
a^{2}+2abi+b^{2}i^{2} &=&3+4i \\
(a^{2}-b^{2})+2abi &=&3+4i
\end{eqnarray*}
i, per tant,
\[
\left\{
\begin{array}{l}
a^{2}-b^{2}=3 \\
2ab=4
\end{array}
\right.
\]
Resolent aquest sistema, s’obtén dues solucions
\[
\left\{
\begin{array}{l}
a=2 \\
b=1
\end{array}
\right. \text{ \ \ \ i \ \ \ }\left\{
\begin{array}{l}
a=-2 \\
b=-1
\end{array}
\right.
\]
Per tant, les arrels quadrades del complex donat són $2+i$ i $
-2-i$.
\end{solucio}

\begin{exercici}
Una arrel cúbica d’un complex és $1-i$. Calcula aquest complex i
les demás arrels.
\end{exercici}

\begin{solucio}
El complex $1-i$ en forma polar és $\left( \sqrt{2}\right) _{-45
^{\circ}
}$. Si $z$ és el complex que busquem, aleshores s’ha de complir
\[
z=\left[ \left( \sqrt{2}\right) _{-45
^{\circ}
}\right] ^{3}
\]
i, per tant,
\[
z=\left( 2\sqrt{2}\right) _{-135
^{\circ}
}=\left( 2\sqrt{2}\right) _{225
^{\circ}
}
\]
és a dir,
\begin{eqnarray*}
z &=&2\sqrt{2}(\cos 225
^{\circ}
+i\sin 225
^{\circ}
) \\
&=&2\sqrt{2}\left( -\frac{\sqrt{2}}{2}-\frac{\sqrt{2}}{2}i\right) \\
&=&-2-2i
\end{eqnarray*}
Com els mòduls de totes les arrels cúbiques de $z$ són iguals
a $\sqrt{2}$ i els arguments de dues arrels consecutives difereixen en
\[
\frac{360
^{\circ}
}{3}=120
^{\circ}
\]
aleshores les restantes arrels són:
\[
\begin{array}{l}
\left( \sqrt{2}\right) _{195
^{\circ}
+120
^{\circ}
}=\left( \sqrt{2}\right) _{75
^{\circ}
} \\
\left( \sqrt{2}\right) _{315
^{\circ}
-120
^{\circ}
}=\left( \sqrt{2}\right) _{195
^{\circ}
}
\end{array}
\]
\end{solucio}

\begin{exercici}
(a) Prova que si $z_{0}$ és una arrel $n$-èsima d’un complex no
nul $z$, aleshores les arrels $n$-èsimes de $z$ són les arrels
$n$-èsimes de la unitat multiplicades per $z_{0}$. (b) Com a aplicació, resol l’exercici anterior.
\end{exercici}

\begin{solucio}
Si $z_{0}$ és una arrel $n$-èsima de $z$, aleshores
\[
z_{0}^{n}=z
\]
Si $z_{1}$ és qualsevol una altra arrel $n$-èsima de $z$, tenim
\[
z_{1}^{n}=z
\]
i, per tant,
\[
\frac{z_{1}^{n}}{z_{0}^{n}}=\left( \frac{z_{1}}{z_{0}}\right) ^{n}=1
\]
Després, el complex $\frac{z_{1}}{z_{0}}$ és una arrel $n$-èsima de
la unitat. Com a conseqüència, qualsevol una altra arrel $n$-èsima de $z$
és el producte de $z_{0}$ per una arrel $n$-èsima de la unitat.

(b) Sabem que una arrel cúbica d’un complex $z$ és $z_{1}=1-i$ i
volem trobar les altres dues arrels. Com les arrels cúbiques de la unitat són
\[
\begin{array}{l}
u_{1}=1_{0
^{\circ}
}=1 \\
u_{2}=1_{120
^{\circ}
}=-\frac{1}{2}+\frac{\sqrt{3}}{2}i \\
u_{3}=1_{240
^{\circ}
}=-\frac{1}{2}-\frac{\sqrt{3}}{2}i
\end{array}
\]
les altres arrels de $z$ són
\[
z_{2}=(1-i)\cdot \left( -\frac{1}{2}+\frac{\sqrt{3}}{2}i\right) =-\frac{1}{2}
+\frac{\sqrt{3}}{2}+\left( \frac{\sqrt{3}}{2}+\frac{1}{2}\right) i
\]
i
\[
z_{3}=(1-i)\cdot \left( -\frac{1}{2}-\frac{\sqrt{3}}{2}i\right) =-\frac{1}{2}
-\frac{\sqrt{3}}{2}+\left( -\frac{\sqrt{3}}{2}+\frac{1}{2}\right) i
\]
\end{solucio}

\begin{exercici}
Resol en $\mathbb{C}$ l’equació següent
\[
(2-i)z^{2}-(3+i)z-2+6i=0
\]
Per què les solucions no són dos complexos conjugats?
\end{exercici}

\begin{solucio}
Al no ser una equació de segon grau amb coeficients reals, no
podem assegurar que les seves solucions siguin dos complexos conjugats.

Resolent l’equació com en el cas del cos $\mathbb{R}$ dels nombres reals, obtenim
\[
\frac{3+i+\sqrt{(3+i)^{2}-4(2-i)(-2+6i)}}{2(2-i)}=\frac{3+i+\sqrt{-50i}}{4-2i
}
\]
Ara bé, com $-50i$ en forma polar és $50_{270
^{\circ}
}$, tenim
\[
\sqrt{50_{270
^{\circ}
}}=\left( \sqrt{50}\right) _{\frac{270
^{\circ}
}{2}+k\cdot \frac{360
^{\circ}
}{2}}=\left\{
\begin{array}{ll}
\left( 5\sqrt{2}\right) _{135
^{\circ}
} & k=0 \\
\left( 5\sqrt{2}\right) _{315
^{\circ}
} & k=1
\end{array}
\right.
\]
i, expressant aquestes arrels en forma binòmica, obtenim
\[
\sqrt{-50i}=\left\{
\begin{array}{l}
-5+5i \\
5-5i
\end{array}
\right.
\]
Per tant,
\begin{eqnarray*}
\frac{3+i+\sqrt{-50i}}{4-2i} &=&\frac{3+i-5+5i}{4-2i} \\
&=&\frac{-2+6i}{4-2i} \\
&=&\frac{(-1+3i)(2+i)}{(2-i)(2+i)} \\
&=&-1+i
\end{eqnarray*}
o bé
\begin{eqnarray*}
\frac{3+i+\sqrt{-50i}}{4-2i} &=&\frac{3+i+5-5i}{4-2i} \\
&=&\frac{8-4i}{4-2i} \\
&=&2
\end{eqnarray*}
Per tant, les solucions de
\[
(2-i)z^{2}-(3+i)z-2+6i=0
\]
són $z=2$ o $z=-1+i$.
\end{solucio}

\begin{exercici}
Resol en $\mathbb{C}$ les equacions següents:

\begin{enumerate}
\item[a)]
\[
z^{4}+64z=0
\]

\item[b)]
\[
z^{6}-28z^{3}+27=0
\]
\end{enumerate}
\end{exercici}

\begin{solucio}
(a)
\[
\begin{array}{lll}
z(z^{3}+64)=0 & \Longrightarrow & \left\{
\begin{array}{l}
z=0 \\
z^{3}+64=0
\end{array}
\right.
\end{array}
\]
Les solucions de l’equació
\[
z^{3}+64=0
\]
són les arrels cúbiques de $-64$. Com $-64$ en forma polar és $
64_{180
^{\circ}
}$, tenim
\[
\sqrt[3]{64_{180
^{\circ}
}}=\left( \sqrt[3]{64}\right) _{\frac{180
^{\circ}
}{3}+k\cdot \frac{360
^{\circ}
}{3}}=\left\{
\begin{array}{ll}
4_{60
^{\circ}
} & k=0 \\
4_{180
^{\circ}
} & k=1 \\
4_{240
^{\circ}
} & k=2
\end{array}
\right.
\]
Passant aquestes arrels a forma binòmica, obtenim
\[
\begin{array}{ll}
\text{Polar} & \text{Binòmica} \\
4_{60
^{\circ}
} & 2+2\sqrt{3}i \\
4_{180
^{\circ}
} & -4 \\
4_{240
^{\circ}
} & 2-2\sqrt{3}i
\end{array}
\]
Per tant, les solucions de l’equació $z^{4}+64z=0$ són $
z_{1}=0$, $z_{2}=2+2\sqrt{3}i$, $z_{3}=-4$ i $z_{4}=2-2\sqrt{3}i$.

(b) Observa que l’equació
\[
z^{6}-28z^{3}+27=0
\]
és biquadrada. Mitjançant el canvi $u=z^{3}$, s’obté l’equació
\[
u^{2}-28u+27=0
\]
Resolent aquesta equació, obtenim
\[
\frac{28+\sqrt{784-108}}{2}=\frac{28\pm 26}{2}=\left\{
\begin{array}{l}
27 \\
1
\end{array}
\right.
\]
Per tant, ara cal resoldre en $\mathbb{C}$ les equacions següents:
\[
z^{3}=27\text{ \ \ \ i \ \ \ }z^{3}=1
\]
Les solucions de la primera són les arrels cúbiques de $27_{0
^{\circ}
}$, és a dir,
\[
\sqrt[3]{27_{0
^{\circ}
}}=\left( \sqrt[3]{27}\right) _{k\cdot \frac{360
^{\circ}
}{3}}=\left\{
\begin{array}{ll}
3_{0
^{\circ}
} & k=0 \\
3_{120
^{\circ}
} & k=1 \\
3_{240
^{\circ}
} & k=2
\end{array}
\right.
\]
i, les solucions de la segona són les arrels cúbiques de $1_{0
^{\circ}
}$, és a dir,
\[
\sqrt[3]{1_{0
^{\circ}
}}=\left( \sqrt[3]{1}\right) _{k\cdot \frac{360
^{\circ}
}{3}}=\left\{
\begin{array}{ll}
1_{0
^{\circ}
} & k=0 \\
1_{120
^{\circ}
} & k=1 \\
1_{240
^{\circ}
} & k=2
\end{array}
\right.
\]
Escribiendo aquestes arrels en forma binòmica, obtenim
\[
\begin{array}{ll}
\text{Polar} & \text{Binòmica} \\
3_{0
^{\circ}
} & 3 \\
3_{120
^{\circ}
} & -\frac{3}{2}+\frac{3\sqrt{3}}{2}i \\
3_{240
^{\circ}
} & -\frac{3}{2}-\frac{3\sqrt{3}}{2}i
\end{array}
\text{ \ \ i \ \ }
\begin{array}{ll}
\text{Polar} & \text{Binòmica} \\
1_{0
^{\circ}
} & 1 \\
1_{120
^{\circ}
} & -\frac{1}{2}+\frac{\sqrt{3}}{2}i \\
1_{240
^{\circ}
} & -\frac{1}{2}-\frac{\sqrt{3}}{2}i
\end{array}
\]
Per tant, les solucions de l’equació
\[
z^{6}-28z^{3}+27=0
\]
són $z_{1}=3$, $z_{2}=1$, $z_{3}=-\frac{3}{2}+\frac{3\sqrt{3}}{2}i$, $z_{4}=-
\frac{1}{2}+\frac{\sqrt{3}}{2}i$, $z_{5}=-\frac{3}{2}-\frac{3\sqrt{3}}{2}i$
i $z_{6}=-\frac{1}{2}-\frac{\sqrt{3}}{2}i$.
\end{solucio}

\begin{exercici}
Troba dos nombres complexos tals que el seu producte és $2i$, i el
quadrat de uno de ellos, dividido pel un altre, és $1/2$.
\end{exercici}

\begin{solucio}
Suposem que $r_{\alpha }$ i $s_{\beta }$ són les formes polars dels complexos que volem trobar. Com el producte de tots dos és $2i$, podem
escriure
\begin{eqnarray*}
r_{\alpha }\cdot s_{\beta } &=&2_{90
^{\circ}
} \\
\left( r\cdot s\right) _{\alpha +\beta } &=&2_{90
^{\circ}
}
\end{eqnarray*}
De aquí, deduïm
\[
\left\{
\begin{array}{l}
r\cdot s=2 \\
\alpha +\beta =90
^{\circ}
+k\cdot 360
^{\circ}
\end{array}
\right.
\]
Per un un altre costat, com el quadrat de uno de ellos, dividido pel un altre és $
1/2 $, podem escriure
\begin{eqnarray*}
\frac{\left( r_{\alpha }\right) ^{2}}{s_{\beta }} &=&\left( \frac{1}{2}
\right) _{0
^{\circ}
} \\
\left( \frac{r^{2}}{s}\right) _{2\alpha -\beta } &=&\left( \frac{1}{2}
\right) _{0
^{\circ}
}
\end{eqnarray*}
De aquí, deduïm
\[
\left\{
\begin{array}{l}
\frac{r^{2}}{s}=\frac{1}{2} \\
2\alpha -\beta =k\cdot 360
^{\circ}
\end{array}
\right.
\]
Per tant, hem obtingut un parell de sistemes:
\[
\left\{
\begin{array}{l}
r\cdot s=2 \\
\frac{r^{2}}{s}=\frac{1}{2}
\end{array}
\right. \text{ \ \ \ i \ \ \ }\left\{
\begin{array}{l}
\alpha +\beta =90
^{\circ}
+k\cdot 360
^{\circ}
\\
2\alpha -\beta =k\cdot 360
^{\circ}
\end{array}
\right.
\]
Del primer, s’obté $r=1$ i $s=2$, i del segon,
\[
\begin{array}{l}
\alpha =30
^{\circ}
+k\cdot 120
^{\circ}
\\
\beta =60
^{\circ}
+k\cdot 120
^{\circ}
\end{array}
\]
Per tant, obtenim tres solucions
\[
\begin{array}{lll}
r_{\alpha } & s_{\beta } &  \\
1_{30
^{\circ}
} & 2_{60
^{\circ}
} & k=0 \\
1_{150
^{\circ}
} & 2_{180
^{\circ}
} & k=1 \\
1_{270
^{\circ}
} & 2_{300
^{\circ}
} & k=2
\end{array}
\]
\end{solucio}

\begin{exercici}
Troba els complexos tals que el quadrat del seu invers és l’oposat del seu conjugat.
\end{exercici}

\begin{solucio}
Si $z$ és un dels complexos que busquem, aleshores ha de complir la
següent condició:
\[
\left( \frac{1}{z}\right) ^{2}=-\overline{z}
\]
Si $z=r_{\alpha }$, aleshores
\[
\begin{array}{lll}
\frac{1}{z}=\left( \frac{1}{r}\right) _{-\alpha } & \Longrightarrow & \left(
\frac{1}{z}\right) ^{2}=\left[ \left( \frac{1}{r}\right) ^{2}
\right] _{-2\alpha }
\end{array}
\]
i
\[
\begin{array}{lll}
\overline{z}=r_{-\alpha } & \Longrightarrow & -\overline{z}=r_{-\alpha +180
^{\circ}
}
\end{array}
\]
Per tant, expressant la condició en forma polar, tenim
\[
\left[ \left( \frac{1}{r}\right) ^{2}\right] _{-2\alpha }=r_{-\alpha +180
^{\circ}
}
\]
Després,
\[
\left\{
\begin{array}{l}
\left( \frac{1}{r}\right) ^{2}=r \\
-2\alpha =-\alpha +180
^{\circ}
+k\cdot 360
^{\circ}
\end{array}
\right.
\]
De la primera equació s’obté $r=1$, i de la segona,
\[
\alpha =180
^{\circ}
\]
Per tant,
\[
z=-1
\]
\end{solucio}

\begin{exercici}
Troba dos complexos del mateix mòdul, els quals arguments sumen $90
^{\circ}
$ i el primer és conjugat del quadrat del segon.
\end{exercici}

\begin{solucio}
Suposem que $r_{\alpha }$ i $s_{\beta }$ són les formes polars dels complexos que volem trobar. Com són del mateix mòdul i els seus arguments sumen $90
^{\circ}
$, podem escriure
\[
\left\{
\begin{array}{l}
r=s \\
\alpha +\beta =90
^{\circ}
\end{array}
\right.
\]
Si el primer és conjugat del quadrat del segon, podem escriure
\[
r_{\alpha }=\left( s^{2}\right) _{-2\beta }
\]
i, per tant, s’ha de complir
\[
\left\{
\begin{array}{l}
r=s^{2} \\
\alpha =-2\beta +k\cdot 360
^{\circ}
\end{array}
\right.
\]
D’aquesta manera, tenim dos sistemes:
\[
\left\{
\begin{array}{l}
r=s \\
r=s^{2}
\end{array}
\right. \text{ \ \ \ i \ \ \ }\left\{
\begin{array}{l}
\alpha +\beta =90
^{\circ}
\\
\alpha =-2\beta +k\cdot 360
^{\circ}
\end{array}
\right.
\]
Del primer s’obté $r=s=1$, i del segon,
\[
\begin{array}{l}
\alpha =180
^{\circ}
\\
\beta =-90
^{\circ}
\end{array}
\]
Per tant, els nombres complexos són $-1$ i $-i$.
\end{solucio}

\begin{exercici}
Donats dos complexos $z_{1}$ i $z_{2}$, interpreta geomètricament la
següent transformació
\[
z~\longmapsto ~z_{1}+z\cdot z_{2}
\]
sent $z\in \mathbb{C}$. Com a aplicació d’aquest resultat, calcula el vèrtex $C$ d’un triangle equilàter $ABC$ si sabem que $
A=(1,1)$ i $B=(4,2)$.
\end{exercici}

\begin{solucio}
Recorda que la transformació $z~\longmapsto ~z+z_{0}$ representa geomètricament una translació, i $z~\longmapsto ~z\cdot z_{0}$, la
composició d’una homotècia i un gir. Com l’operació $
z_{1}+z\cdot z_{2}$ és una combinació de les dues anteriors, la
transformació $z~\longmapsto ~z_{1}+z\cdot z_{2}$ representa geomètricament la composició d’una homotècia i un gir seguida d’una
translació. Per exemple, en la figura següent\begin{center}
\includegraphics[width=1.58in,height=1.7677in,keepaspectratio]{com-15.png}
\end{center}es compleix
\[
z_{C}=z_{A}+(z_{B}-z_{A})\cdot r_{\theta }
\]
sent
\[
r=\left\vert \frac{z_{C}-z_{A}}{z_{B}-z_{A}}\right\vert
\]

Considerem els complexos $z_{A}$ i $z_{B}$ els quals afixos són els punts $
A=(1,1)$ i $B=(4,2)$. Aleshores, és claro que el complex $z_{B}-z_{A}$ té
per mòdul la longitud del costat del triangle equilàter. A més, com tots els angles d’aquest triangle mesuren $60
^{\circ}
$, es compleix
\[
z_{C}=z_{A}+(z_{B}-z_{A})\cdot 1_{60
^{\circ}
}
\]
o bé,
\[
z_{C}=z_{A}+(z_{B}-z_{A})\cdot 1_{-60
^{\circ}
}
\]
\begin{center}
\includegraphics[width=1.8507in,height=1.8715in,keepaspectratio]{com-8.png}
\end{center}Per tant, prenent la primera equació, obtenim
\begin{eqnarray*}
z_{C} &=&1+i+(3+i)\cdot (\cos 60
^{\circ}
+i\sin 60
^{\circ}
) \\
&=&1+i+(3+i)\cdot \left( \frac{1}{2}+\frac{\sqrt{3}}{2}i\right) \\
&=&1+i+\frac{3}{2}+\frac{1}{2}i+\frac{3\sqrt{3}}{2}i+\frac{\sqrt{3}}{2}i^{2}
\\
&=&\frac{5-\sqrt{3}}{2}+\frac{3+3\sqrt{3}}{2}i
\end{eqnarray*}
i, com a conseqüència, les coordenades del punt $C$ són
\[
C=\left( \frac{5-\sqrt{3}}{2},\frac{3+3\sqrt{3}}{2}\right)
\]
De la mateixa manera, prenent la segona equació, s’obté l’altra solució
\[
C^{\prime }=\left( \frac{5+\sqrt{3}}{2},\frac{3-3\sqrt{3}}{2}\right)
\]
\end{solucio}

\begin{exercici}
Determina les coordenades dels vèrtexs $C$ i $D$ del quadrat $ABCD$,
si $A=(1,3)$ i $B=(5,2)$.
\end{exercici}

\begin{solucio}
Suposem que $z_{A}$, $z_{B}$, $z_{C}$ i $z_{D}$ són els complexos els quals
afixos són els vèrtexs del quadrat. Com que es tracta d’un quadrat, els
complexos $z_{B}-z_{A}$ i $z_{D}-z_{A}$ tenen el mateix mòdul i l’angle que determinen els seus vectors associats és recte. Per tant, es compleix
\[
z_{D}=z_{A}+(z_{B}-z_{A})\cdot 1_{90
^{\circ}
}
\]
o bé,
\[
z_{D}=z_{A}+(z_{B}-z_{A})\cdot 1_{-90
^{\circ}
}
\]
\begin{center}
\includegraphics[width=1.58in,height=1.7573in,keepaspectratio]{com-14.png}
\end{center}

Prenent la primera equació, obtenim
\[
z_{D}=1+3i+(4-i)\cdot i=2+7i
\]
i, per tant, el punt $D=(2,7)$. Per determinar $C$, observem que també
es compleix
\[
z_{C}=z_{B}+(z_{A}-z_{B})\cdot 1_{-90
^{\circ}
}
\]
Després,
\[
z_{C}=5+2i+(-4+i)\cdot (-i)=6+6i
\]
i, per tant, el punt $C=(6,6)$.

De la mateixa manera, prenent la segona equació, s’obté $D=(0,-1)$ i
después $C=(4,-2)$.
\end{solucio}

\begin{exercici}
Donats els nombres complexos $z_{1}=-1+i$ i $z=-2+3i$, calcula dos nombres complexos $z_{2},z_{3}$ tals que els afixos de $z_{1},z_{2}$ i $z_{3}$
formen un triangle equilàter de centre l’afix de $z$.
\end{exercici}

\begin{solucio}
Sabem que els afixos de les arrels cúbiques d’un nombre
complex $w$ són els vèrtexs d’un triangle equilàter inscrit en
una circumferència de centre l’origen i radi $\sqrt[3]{\left\vert
w\right\vert }$. Com que el nostre triangle ha d’estar centrat en el
afix de $z$, haurem d’efectuar una translació de vector associat a $
z$. D’aquesta manera, tenim
\[
z_{i}=z+\sqrt[3]{w}~\ \ (i=1,2,3)
\]
Com una de les arrels cúbiques de $w$ és $z_{1}-z$, podem obtenir
les arrels cúbiques de $w$ multiplicant $z_{1}-z$ per les arrels cúbiques de la unitat (vegeu l’exercici anterior). Les arrels cúbiques de la unitat són
\[
\begin{array}{l}
u_{1}=1 \\
u_{2}=\cos \frac{2\pi }{3}+i\sin \frac{2\pi }{3}=-\frac{1}{2}+i\frac{\sqrt{3}
}{2} \\
u_{3}=\cos \frac{4\pi }{3}+i\sin \frac{4\pi }{3}=-\frac{1}{2}-i\frac{\sqrt{3}
}{2}
\end{array}
\]
Per tant, les arrels cúbiques de $w$ són
\[
\begin{array}{l}
w_{1}=1-2i \\
w_{2}=(1-2i)\left( -\frac{1}{2}+i\frac{\sqrt{3}}{2}\right) =-\frac{1}{2}+
\sqrt{3}+i\left( \frac{\sqrt{3}}{2}+1\right) \\
w_{3}=(1-2i)\left( -\frac{1}{2}-i\frac{\sqrt{3}}{2}\right) =-\frac{1}{2}-
\sqrt{3}+i\left( -\frac{\sqrt{3}}{2}+1\right)
\end{array}
\]
i, en conseqüència, els nombres complexos que busquem són
\[
\begin{array}{l}
z_{1}=z+w_{1}=-1+i \\
z_{2}=z+w_{2}=-\frac{5}{2}+\sqrt{3}+i\left( \frac{\sqrt{3}}{2}+4\right) \\
z_{3}=z+w_{3}=-\frac{5}{2}-\sqrt{3}+i\left( -\frac{\sqrt{3}}{2}+4\right)
\end{array}
\]
\end{solucio}
